Neither page fault tracepoint is exported to modules, so the module cannot
call \path{register_trace_page_fault_user()}. Instead,
\path{for_each_kernel_tracepoint()} visits every tracepoint built into
the kernel, \texttt{find\_tracepoint()} keeps the two whose names match, and
\texttt{tracepoint\_probe\_register()} attaches a probe to each. Attaching the
first probe also runs the tracepoints' registration hook,
\texttt{trace\_pagefault\_reg()}, which turns on the static key that the fault
handler tests; until then the handler skips the tracing code altogether.

A probe runs on every page fault in the system with preemption disabled, so
it does one comparison: the thread group of the faulting task against the
monitored one. It compares \texttt{struct pid} pointers rather than numbers,
so a reused process ID cannot be mistaken for the monitored process, and it
counts the faults of every thread of the process. To change the target, the
module first clears it and calls \texttt{synchronize\_rcu()}, which waits for
any probe still running; only then does it reset the counters. On unloading,
\texttt{tracepoint\_synchronize\_unregister()} ensures no probe is still
executing module code that is about to be freed.
